\section{Results}
\label{sec:results}

\subsection{M0 Baseline Results (RQ1)}

\begin{table}[h]
\centering
\caption{M0 (Full KV) per-task F1 on LongBench 4-task QA.}
\label{tab:m0}
\begin{tabular}{ll}
\toprule
Task & M0 F1\% (raw) \\
\midrule
NarrativeQA & 9\% (0.09) \\
HotpotQA & 9\% (0.09) \\
2WikiMQA & 11\% (0.11) \\
MuSiQue & 6\% (0.06) \\
\midrule
\textbf{Macro-average} & \textbf{8.75\% (0.0875)} \\
\bottomrule
\end{tabular}
\end{table}

M0 produces coherent, non-degenerate English-language text on all 4 tasks. The macro-average of 8.75\% is consistent with LLaMA-2-7B-chat at 4K context truncation---most relevant context for these multi-hop tasks exceeds the 4K window, so below-state-of-the-art F1 is expected and does not indicate a pipeline error. F1 computation matches the LongBench standard (SQuAD-style normalization).

\subsection{KV Shape Verification (RQ2)}

The shape log confirms that M1 eviction reduces KV dimensions correctly across all 32 transformer layers:

\begin{verbatim}
Layer  0: [1, 32, 4096, 128] -> [1, 32, 2048, 128]  OK
Layer  1: [1, 32, 4096, 128] -> [1, 32, 2048, 128]  OK
  ...
Layer 31: [1, 32, 4096, 128] -> [1, 32, 2048, 128]  OK
32/32 layers confirmed at 2048/4096 positions per head.
\end{verbatim}

\subsection{M1 Generation Quality (RQ3)}

\begin{table}[h]
\centering
\caption{M1 (Prefill-Observation) vs.\ M0 F1.}
\label{tab:m1_vs_m0}
\begin{tabular}{llll}
\toprule
Task & M0 F1\% (raw) & M1 F1\% (raw) & Status \\
\midrule
NarrativeQA & 9\% (0.09) & 0\% (0.00) & Degenerate \\
HotpotQA    & 9\% (0.09) & 0\% (0.00) & Degenerate \\
2WikiMQA    & 11\% (0.11) & 0\% (0.00) & Degenerate \\
MuSiQue     & --- & --- & Not executed \\
\bottomrule
\end{tabular}
\end{table}

Despite passing shape verification, M1 produces F1=0\% on all 3 completed tasks. F1=0\% requires zero token overlap with the reference answer---indicating complete generation failure, not suboptimal performance.

\subsection{Root Cause Analysis}

\begin{table}[h]
\centering
\caption{Diagnostic chain for M1 failure.}
\label{tab:diagnosis}
\begin{tabular}{llll}
\toprule
Pipeline Step & Expected & Observed & Status \\
\midrule
M0 end-to-end & Coherent generation, F1>0\% & F1=8.75\%, coherent & PASS \\
Score fn (M1) & Valid importance scores & Non-trivial tensor & PASS \\
KV shape eviction & 2048/4096 per head/layer & Confirmed, 32 layers & PASS \\
DynamicCache recon. & Valid cache state & Incorrect internal state & FAIL \\
M1 generation & F1>0\% & F1=0\% (degenerate) & FAIL \\
\bottomrule
\end{tabular}
\end{table}

The causal attribution is clear: all steps before cache reconstruction pass; all steps after it fail. M0's working condition differs from M1's broken condition by exactly one step---the cache reconstruction.

\subsection{Summary}

\begin{table}[h]
\centering
\caption{Research question findings summary.}
\label{tab:rq_summary}
\begin{tabular}{ll}
\toprule
RQ & Finding \\
\midrule
RQ1: M0 correct? & YES --- macro-F1=8.75\%, coherent generation \\
RQ2: Shape eviction correct? & YES --- 2048/4096 per head, 32/32 layers \\
RQ3: M1 non-degenerate? & NO --- F1=0\%, DynamicCache failure \\
H-E1 gate evaluable? & NO --- M2 not executed; M1 degenerate \\
\bottomrule
\end{tabular}
\end{table}
